\documentclass[../main]{subfiles}
\begin{document}
\ifSubfilesClassLoaded{\setcounter{page}{47}}{}
\section{Communism and the Personality}
\label{sec:08_communism_personality}

The problem of the personality is one of the central problems of communism. Under communism, the significance of the personality rises to heights never before seen. “In place of the old bourgeois society, with its classes and class antagonisms, we shall have an association, in which the free development of each is the condition for the free development of all,” wrote Marx and Engels in the “Manifesto of the Communist Party”.

All of preceding history must be understood not simply as a particular sequence of events, but as the process of the becoming of the human essence. It then becomes clear that this is, at the same time, the process of the becoming of a particular type of personality, one that corresponds to the human essence. What should this personality be like? This question will be resolved each time practically and in concrete-historical terms, but in any case it is a personality that consciously creates history within a collective of other creative personalities like itself. In a free association of people, their entire life depends on creative activity, which is the highest expression of the personality that carries it out.

The personality is a social phenomenon from beginning to end. An individual is not born a personality but becomes one, entering into social relations as an actively engaged subject. Initially, human life is given to the newborn from outside—from the social space in which they live, with all the variations of the relations of human being to human being and of human being to nature. Only by assimilating, that is, by making their own, the forms of activity—or, what amounts to the same thing, the forms of culture, which are in no way inherent in the individual from the start—do they become a personality. Everyone who comes into the world finds these forms already in place. For the biological organism they are external, but only for the organism, not for the personality. The world of activity, as the world of human culture, is the real world of the personality. In communist society, everyone must make this world their own, and that means learning to change it.

The communist personality will, in its practice, actively create social relations rather than simply being their passive product. The problem is to create conditions in which the personality, in its individual development, would travel in the shortest possible time the path of society’s becoming and arrive at the forefront of social life. Each time, one must seek and find concrete forms of collectivity that would bind the personality to society in such a way that it could not help being an active member of it and living the full life of society. Since society is constantly developing, such forms of “joint-separate activity” cannot be found once and for all. The management of social development under communism consists mainly in learning to change, in good time, the forms of connection between the human being and society so that they compel people to act humanly in any new concrete-historical conditions.

Communists are accused of wanting to standardise the personality, to make people “all the same”. In making this accusation, the supporters of capitalism are in fact giving away the secret of modern capitalist society. Under capitalism, individuals by and large turn out so standardised that their individuality is defined by mere trifles. This bad individuality has nothing to do with true individuality, understood as the individual expression of the universal human culture a person has assimilated. Capitalism draws people into standard practices of production and consumption and standardises their behaviour both at work and outside it. People are already “all the same”. Communism, by contrast, presupposes the creation of a social space for the development of individuality.

Communists see the need to create a society of universally developed creative personalities. In such a society, the individuality of each person will be the active expression and embodiment of social development. In what sense, then, can we speak of standards here?

Clearly, a person cannot be a universal creator without having assimilated the minimum of human culture necessary for this, and so we must speak of a minimum of conditions that society must provide to each of its members so that they can be a universal creator; there must be a minimum standard for these conditions. What does this minimum consist of?

It includes: a certain level of development of the healthcare system, one aimed specifically at preserving health and creating the conditions for a healthy life for every person, rather than merely treating illness; a quality and level of education that is compulsory for everyone; a system of public catering and, more broadly, of all publicly accessible services capable of abolishing everyday life; and provision for childhood and old age. In one of his works, Lenin wrote that socialism begins with a glass of milk for every child. Obviously, the presence of that glass does not in itself mean that society is developing in the direction of communism, but without children being well fed, it certainly cannot develop in that direction. Nor can it develop without everything listed above (the list, of course, is not complete. It should be drawn up by people with expertise in various fields, and it will expand as society develops).

This minimum of humane treatment of one person by another cannot be fixed once and for all. It changes depending on the conditions in which the transition takes place, that is, on the level of development of human culture. Today’s minimum, the one that would give society the impetus to develop towards communism, must necessarily include higher education encompassing not only computer literacy at the level of programming but also the scientific foundations for understanding all production processes—polytechnism (for everyone!). Ensuring access to all modern means of communication is likewise essential.

It is important to note that even the most developed capitalist countries cannot boast of universal literacy among their entire population or of universally accessible medical care.

The necessary minimum of humaneness in how people are treated is, as already noted, a historical matter. Yet this minimum must also include the need for its own constant expansion. For a communist society, in which the free development of each is the condition for the free development of all, to exist and develop, this minimum must continually incorporate ever newer, most advanced achievements of humankind. Communist society will consist of intelligent, healthy and talented people, who, in the final analysis, are the only ones capable of organising the social reproduction of their lives on rational foundations.
\end{document}
